\section{Future Work}
\label{sec6.3: future work}

The challenges identified in \ref{sec6.2: limitations} point toward several avenues for future research:

\begin{itemize}
    \item \textbf{Reducing Prediction Set Sizes Directly:} Beyond the Gram type filter we already used, we can use a second
          coarser filtering stage, for example something based on phage morphology or a nearest neighbor prefilter could narrow 
          the host level candidate space even further before the conformal envelope test is applied. 

    \item \textbf{Pooling for Low Sample Hosts:} Rather than calibrating all of the $54$ host envelopes 
           independently, a hierarchical shrinkage scheme could pool information across related hosts. Shrinking quantile 
           estimates for data scarce genera ($n_h < 10$) toward family or Gram level baselines would stabilize coverage 
           without requiring additional data.

%     \item \textbf{Dimensionality Reduction:} To resolve high dimensional sparsity in multi 
%            dimensional geometries ($K_h \approx 40$), future work should explore learning host specific linear or non linear 
%            projections $\Phi_h: \mathbb{R}^{K_h} \to \mathbb{R}^{d}$ ($d \in [2, 5]$) designed to maximize infector/non 
%            infector separation. 
        %    Learning a Mahalanobis distance metric may help us preserve local support for the Radial envelope in higher 
        %    dimensions and reduce reliance on the marginal fallback.

    \item \textbf{Specificity Control:} The impostor experiment (\ref{sec5.3.7: host ablations}) shows 
          that naively incorporating non infector scores collapses the prediction sets to near total abstention. Future work 
          may investigate two sided nonconformity functions that jointly penalize false negatives and false positives, 
          this may bound the false discovery rate and yet avoid the coverage loss.
    
%     \item \textbf{Handling of Missing Scores:} Instead of defaulting to the marginal quantile when 
%           functional categories are absent, future envelope constructions could condition fallbacks on the observed 
%           missingness patterns. This may be done by observing the sparsity pattern, which itself may be carrying some 
%           information about the host.
     \item \textbf{Missingness Aware Envelope:} The current Strip construction handles structural missingness by removing
            unobserved dimensions from the pairwise envelope. Future work could investigate whether the observed missingness pattern 
            itself contains some information about the host and if it could be used to get a missingness aware conditional envelope 
            construction.

      \item \textbf{Dimensionality Reduction} As demonstrated by our analysis of core functions (Section~\ref{sec5.3.6:core-set}), 
             most discriminative information lives in a low dimensional subspace. In higher dimensions ($K_h 
           \approx 40$), sample sparsity forces directional cones and grid bins to fall back onto conservative marginal 
           bounds. Future work should explore linear dimensionality reduction such as class specific Principal Component 
           Analysis (PCA) to extract dominant functional variance directions. Supervised linear/non-linear projections 
           $\Phi_h: \mathbb{R}^{K_h} \to \mathbb{R}^d$ ($d \in [2, 5]$) optimized to maximize infector/non infector 
           separation. Calibrating Radial and Strip envelopes within these compact, orthogonal PCA subspaces would 
           eliminate the curse of dimensionality while preserving directional sensitivity/

      % \item \textbf{Per Direction Quantile Level:} In this work the strip envelope calibrates each individual bound (each $(j,i)$ 
      %        constraint) at the $1 \alpha$ quantile level, and the final envelope is the intersection of all of these bounds across 
      %        every pair of dimensions. We then calibrate every direction at $1-\alpha$ and then take the intersection of then which 
      %        means that we throw away $\alpha$ fraction of the data from each direction. So, when the number of dimensions grow, the 
      %        the combined envelope ends up much smaller than intended. This makes the shape discovery envelope overly tight, and forces 
      %        the size scaling factor $\hat t$ to grow much larger than it otherwise would to restore the target coverage. A simple fix 
      %        would be to calibrate each individual direction at a stricter quantile level of $1-\alpha/d$ (where $d$ is the number of 
      %        dimensions), so that each direction only removes just its share of the data and the intersection ends up close to the 
      %        intended $1-\alpha$ level overall. This should also produce a $\hat t$ much closer to $1$ and a better shaped, less 
      %        conservative envelope.
      \item \textbf{Per Direction Quantile Level:} For the conservativeness of strip envelopes, a next step is to calibrate 
             each individual direction at a stricter quantile level of $1-\alpha/d$ (where $d$ is the number of dimensions). This 
             would ensure each that direction only removes a small of data. This would allow the final intersection to match the 
             target $1-\alpha$ coverage, and produce a $\hat{t}$ that is closer to $1$ and is less conservative.
\end{itemize}









